% 层宽 3,4,3,2。浅色神经元、细连接与少量方向提示构成统一视觉语法。
\begin{tikzpicture}[x=1mm,y=1mm,font=\figurefont\footnotesize]
  \foreach \i/\y in {1/28,2/14,3/0}
    \coordinate (x\i) at (12,\y);
  \foreach \i/\y in {1/35,2/21,3/7,4/-7}
    \coordinate (h\i) at (52,\y);
  \foreach \i/\y in {1/28,2/14,3/0}
    \coordinate (k\i) at (92,\y);
  \foreach \i/\y in {1/21,2/7}
    \coordinate (a\i) at (132,\y);
  \foreach \i in {1,2,3} \foreach \j in {1,2,3,4}
    \draw[nn link,draw=NNInput!58,shorten <=4.4mm,shorten >=4.4mm] (x\i)--(h\j);
  \foreach \i in {1,2,3,4} \foreach \j in {1,2,3}
    \draw[nn link,draw=NNHidden!58,shorten <=4.4mm,shorten >=4.4mm] (h\i)--(k\j);
  \foreach \i in {1,2,3} \foreach \j in {1,2}
    \draw[nn link,draw=NNOutput!58,shorten <=4.4mm,shorten >=4.4mm] (k\i)--(a\j);
  \foreach \i in {1,2,3}
    \node[nn input] at (x\i) {$x_{\i}$};
  \foreach \i in {1,2,3,4}
    \node[nn hidden] at (h\i) {$h_{\i}$};
  \foreach \i in {1,2,3}
    \node[nn hidden] at (k\i) {$h_{\i}$};
  \foreach \i in {1,2}
    \node[nn output] at (a\i) {$a_{\i}$};
  \foreach \x/\title/\symbol in {
    12/输入层/\bm{x},52/隐藏层1/\bm{h}^{(1)},
    92/隐藏层2/\bm{h}^{(2)},132/输出层/\bm{a}}{
    \node[nn heading] at (\x,55) {\title};
    \node[text=BookMuted] at (\x,47) {$\symbol$};
  }
  \draw[nn flow] (26,47)--(38,47);
  \draw[nn flow,draw=NNHidden] (66,47)--(78,47);
  \draw[nn flow,draw=NNOutput] (106,47)--(118,47);
  \node[nn annotation] at (72,-19)
    {$\bm{z}^{(\ell)}=\bm{W}_{\ell}\bm{h}^{(\ell-1)}+\bm{b}_{\ell}$
     \qquad $\bm{h}^{(\ell)}=\sigma_{\ell}(\bm{z}^{(\ell)})$};
  \node[nn annotation] at (132,-7) {$\widehat{\bm{y}}=g(\bm{a})$};
  % 无底色、无外框图例，沿用函数图的短样线与文本排列。
  \foreach \x/\col/\txt in {5/NNInput/输入坐标,57/NNHidden/隐藏响应,109/NNOutput/原始得分}{
    \draw[\col,line width=.9pt] (\x,-32)--++(10,0);
    \node[circle,draw=\col,fill=\col!16,line width=1.1pt,
      minimum size=3.6mm,inner sep=0pt] at (\x+5,-32) {};
    \node[anchor=west] at (\x+13,-32) {\txt};
  }
\end{tikzpicture}
